\begin{theorem}[Tame odd preparation]
\label{mod:cases-tameoddpreparation}
Select the minimal norm-character-orbit representative, route conductor zero,
one, high even, and high odd through the completed tame-odd APIs, supply the
prepared ramified Hasse input only in the high odd branch, bridge the full
product identity (including the trivial norm character), and transport it back
to the original character by the proved orbit invariances.

For high conductor, use the actual norm and trace pullbacks, decompose each
actual conductor by quotient and remainder modulo two, choose admissible
denominators, and choose the stationary representative as a preimage of the
stationary coefficient quotient class.  In the odd branch, use the lower and
upper source-uniformizer coordinates and the resulting prepared ramified Hasse
comparison.  The final theorem is dispatch-ready and concludes the literal
\texttt{FirstMainIdentity}.
\LeanFile{LanglandsFirstMainLemma/Cases/TameOddPreparation.lean}
\lean{LanglandsFirstMainLemma.firstMain_tameOdd_of_isDeltaFinite}
\leanok
\SourceText{Proposition prop:tame-odd-new and Section sec:first-main-completion.}
\uses{mod:ramification-primecyclicpreparation,mod:parameters-minimalorbitstationary,mod:parameters-ramifiedhassepreparation,mod:cases-tameodd,mod:delta-errorterm}
\FormalizationContract
No new Hasse calculation belongs here, and no stationary data are introduced
in conductors zero or one.  The tame hypothesis is used only to identify the
break with zero.  If twisting lowers a conductor, build the stationary class
at the new depth; never reuse a cancelled sum from the old layer.  Transport
the elementary stationary value and residual Hasse phase together.
\end{theorem}
